\section*{Lectures on Lie Groups and Lie Algebras – 13,14}

(The material of this lecture appears in much greater detail in
Chapter 5 of Serre’s book \emph{Lie Algebras and Lie Groups}.)

\section*{Nilpotent and Solvable Lie Algebras}

\textbf{Definition 1.}
A Lie algebra $\mathfrak g$ is called \emph{nilpotent} if there exists
$N \in \mathbb{N}$ such that for any $x_1,\dots,x_N \in \mathfrak g$,

\[
[x_1,[x_2,\dots,[x_{N-1},x_N]\dots]]=0.
\]

Equivalently,
\[
[\mathfrak g,[\mathfrak g,\dots[\mathfrak g,\mathfrak g]\dots]]=0.
\]

Nilpotent Lie groups are defined analogously.

Nilpotency is equivalent to the existence of a basis in which
all operators $\mathrm{ad}\,x$ are strictly upper triangular matrices.

\textbf{Definition 2.}
The \emph{derived series} of $\mathfrak g$ is defined inductively by

\[
D^0\mathfrak g=\mathfrak g,
\qquad
D^k\mathfrak g=[D^{k-1}\mathfrak g,D^{k-1}\mathfrak g].
\]

The Lie algebra $\mathfrak g$ is \emph{solvable} if
$D^N\mathfrak g=0$ for some $N$.

\textbf{Proposition 1.}
Every nilpotent Lie algebra is solvable.

\textbf{Examples.}

1. $\mathfrak n_+(n)$ (strictly upper triangular matrices) is nilpotent.

2. $\mathfrak b_+(n)$ (upper triangular matrices) is solvable but not
nilpotent for $n>1$.

\textbf{Proposition 2.}
$\mathfrak g$ is solvable iff there exists a chain of ideals

\[
0=\mathfrak g_0\subset\mathfrak g_1\subset\dots\subset\mathfrak g_N=\mathfrak g
\]

such that each quotient $\mathfrak g_{i+1}/\mathfrak g_i$ is abelian.

\textbf{Proposition 3.}
Subalgebras and quotients of solvable (resp. nilpotent) Lie algebras
are solvable (resp. nilpotent).
An extension of a solvable Lie algebra by a solvable one is solvable.

\textbf{Proposition 4.}
If $G$ is connected and simply connected with solvable Lie algebra
$\mathfrak g$, then

\[
\exp:\mathfrak g\to G
\]

is a diffeomorphism.
If $\mathfrak g$ is nilpotent, the group multiplication
is given by polynomial formulas.

\textbf{Corollary 1.}
$G$ is solvable (resp. nilpotent)
iff $\mathfrak g=\mathrm{Lie}\,G$ is solvable (resp. nilpotent).

\section*{Engel's Theorem}

\textbf{Theorem 1 (Engel).}
Let $\mathfrak g\subset\mathfrak{gl}(V)$ consist of nilpotent operators.
Then there exists a basis of $V$ in which all elements of $\mathfrak g$
are strictly upper triangular.

In particular, $\mathfrak g$ is nilpotent.

\textbf{Corollary 2.}
$\mathfrak g$ is nilpotent iff $\mathrm{ad}\,x$ is nilpotent
for every $x\in\mathfrak g$.

\section*{Lie's Theorem}

\textbf{Theorem 2 (Lie).}
Let $\mathfrak g\subset\mathfrak{gl}(V)$ be solvable over $\mathbb C$.
Then there exists a basis of $V$ in which all elements
of $\mathfrak g$ are upper triangular.

(Note: false over $\mathbb R$.)

\textbf{Corollary 3.}
The commutator $[\mathfrak g,\mathfrak g]$
of a solvable Lie algebra is nilpotent.

\section*{Cartan's Criterion and Semisimple Lie Algebras}

\textbf{Theorem 3 (Cartan's Criterion).}
Let $\mathfrak g\subset\mathfrak{gl}(V)$ be such that

\[
\mathrm{tr}(xy)=0
\quad\text{for all } x\in\mathfrak g,\ y\in[\mathfrak g,\mathfrak g].
\]

Then $\mathfrak g$ is solvable.

\textbf{Corollary 4.}
If the Killing form of $\mathfrak g$ is identically zero,
then $\mathfrak g$ is solvable.

\textbf{Definition 3.}
A Lie algebra $\mathfrak g$ is \emph{semisimple}
if it contains no nontrivial solvable ideals.

\textbf{Examples.}

1. $\mathfrak{sl}_n(\mathbb C)$ is simple (hence semisimple).

2. The Lie algebra of affine transformations of the plane
is neither semisimple nor solvable.

\textbf{Corollary 5.}
$\mathfrak g$ is semisimple iff its Killing form is nondegenerate.

\textbf{Corollary 6.}
A semisimple Lie algebra decomposes as a direct sum
of nonabelian simple Lie algebras.

\textbf{Proposition 5.}
Every complex Lie algebra has a maximal solvable ideal,
called the \emph{solvable radical}.
The quotient by the radical is semisimple.